%% =============================================================================
%% BEATRIX - Innosuisse Innovation Cheque Antragstext
%% Version 5.2 - Basiert auf Originaldokument (10.2.2026) mit:
%%   (a) BCM/EBF/BEATRIX 3-Ebenen-Abgrenzung (Ernst Fehr Korrektur)
%%   (b) "Neue Absätze" (Harald Gall Feedback) als Basis
%%   (c) Formaler, professioneller Ton des Originals beibehalten
%%   (d) 7-Kapitel-Struktur des Originals beibehalten
%%   (e) Anwendungsbeispiel 3 (M&A) überarbeitet und gestrafft
%% Datum: 16. Februar 2026
%% Terminologie:
%%   BCM  = Beratungsmodell (30+ Jahre Praxis, mentales Rahmenwerk)
%%   EBF  = Wissenschaftliche Formalisierung (enthält BCM, geht darüber hinaus)
%%   BEATRIX = Technologiegestützte Infrastruktur (operationalisiert EBF)
%%   Beziehung: BCM ⊂ EBF ← BEATRIX
%% =============================================================================

\documentclass[11pt,a4paper]{article}

%% Packages
\usepackage[utf8]{inputenc}
\usepackage[T1]{fontenc}
\usepackage[ngerman]{babel}
\usepackage{geometry}
\usepackage{booktabs}
\usepackage{longtable}
\usepackage{array}
\usepackage{xcolor}
\usepackage{hyperref}
\usepackage{enumitem}
\usepackage{titlesec}
\usepackage{fancyhdr}
\usepackage{parskip}
\usepackage{microtype}
\usepackage{graphicx}
\usepackage{amssymb}

%% Page geometry
\geometry{
    left=2.5cm,
    right=2.5cm,
    top=2.5cm,
    bottom=2.5cm
}

%% Colors
\definecolor{fehrgray}{RGB}{80,80,80}
\definecolor{fehrblue}{RGB}{0,51,102}

%% Hyperref setup
\hypersetup{
    colorlinks=true,
    linkcolor=fehrblue,
    citecolor=fehrblue,
    urlcolor=fehrblue
}

%% Section formatting
\titleformat{\section}
    {\Large\bfseries\color{fehrblue}}
    {\thesection}{1em}{}
\titleformat{\subsection}
    {\large\bfseries\color{fehrgray}}
    {\thesubsection}{1em}{}
\titleformat{\subsubsection}
    {\normalsize\bfseries}
    {\thesubsubsection}{1em}{}

%% Header/Footer
\pagestyle{fancy}
\fancyhf{}
\fancyhead[L]{\small\textcolor{fehrgray}{BEATRIX -- Evidence Foundation -- FehrAdvice \& Partners AG}}
\fancyhead[R]{\small\textcolor{fehrgray}{SSBM}}
\fancyfoot[C]{\thepage}
\fancyfoot[R]{\small Version: 5.2}
\renewcommand{\headrulewidth}{0.4pt}

%% Table column types
\newcolumntype{L}[1]{>{\raggedright\arraybackslash}p{#1}}
\newcolumntype{C}[1]{>{\centering\arraybackslash}p{#1}}

\begin{document}

%% =============================================================================
%% TITELSEITE
%% =============================================================================

\begin{titlepage}
\centering
\vspace*{2cm}

{\Huge\bfseries\color{fehrblue} BEATRIX}\\[0.5cm]
{\Large Behavioral Economics AI Technology for\\Research and Implementation eXcellence}\\[2cm]

{\LARGE Validierung einer KI-gestützten\\Beratungs- und Prognoseinfrastruktur\\zur Verbesserung komplexer\\organisationaler Entscheidungen}\\[1.5cm]

{\large Version 5.2}\\[2cm]

\rule{0.8\textwidth}{0.5pt}\\[1.5cm]

\begin{tabular}{ll}
\textbf{Förderprogramm:} & Innosuisse Innovation Cheque \\[0.3cm]
\textbf{Fördersumme:} & CHF 15'000 \\[0.3cm]
\textbf{Fachgebiet:} & Social Sciences \& Business Management (SSBM) \\[0.3cm]
\textbf{Wirtschaftspartner:} & FehrAdvice \& Partners AG \\[0.3cm]
\textbf{Forschungspartner:} & Universität Zürich \\[0.3cm]
& Prof. Harald Gall \quad Institut für Informatik, UZH \\
& Prof. Johannes Luger \quad Institut für Betriebswirtschaft, UZH \\
& Prof. Ernst Fehr \quad Scientific Advisor \\
\end{tabular}

\vfill

\fbox{\parbox{0.9\textwidth}{\centering\large\bfseries
\textit{``Die Technologie ermöglicht die Innovation, aber die Transformation der\\Beratungspraxis ist das Produkt.''}
}}

\vfill

{\large Datum: 14. Februar 2026}

\end{titlepage}

%% =============================================================================
%% INHALTSVERZEICHNIS
%% =============================================================================

\tableofcontents
\newpage

%% =============================================================================
%% 1. EXECUTIVE SUMMARY
%% =============================================================================

\section{Executive Summary}

BEATRIX ist eine KI-gestützte Entscheidungs-, Prognose- und Beratungsinfrastruktur, die auf den Ergebnissen und der Nutzbarmachung von über drei Jahrzehnten verhaltensökonomischer Forschung beruht. Sie adressiert ein zentrales strukturelles Problem organisationaler Entscheidungsprozesse: die begrenzte Fähigkeit, komplexe Entscheidungszusammenhänge systematisch, konsistent und vergleichbar zu analysieren.

Im Unterschied zu bestehenden Entscheidungs- und Beratungsansätzen berücksichtigt BEATRIX bei jeder Analyse das gesamte im \textbf{EBF} (Evidence-Based Framework) formalisierte verhaltensökonomische Wissen, einschliesslich des theoretisch verfügbaren Wissens, der empirisch belegten Wechselwirkungen sowie der Heterogenität menschlichen Verhaltens. Diese umfassende Verarbeitung übersteigt die kognitiven Möglichkeiten einzelner Personen und Teams -- selbst hochqualifizierter Expertinnen und Experten -- und verschafft BEATRIX eine strukturelle Überlegenheit gegenüber der bisherigen Praxis.

\vspace{0.5cm}
\noindent\textbf{Begriffliche Einordnung: BCM, EBF und BEATRIX.} Drei Begriffe strukturieren diesen Antrag:

\begin{itemize}[nosep]
    \item \textbf{BCM} (Behavioral Change Model): Das Entscheidungs- und Beratungsmodell, das FehrAdvice \& Partners über mehr als drei Jahrzehnte in enger Verbindung mit der Forschung von Prof. Ernst Fehr entwickelt hat. Es wurde von allen Berater:innen als mentales Rahmenwerk eingesetzt -- mächtig, aber begrenzt durch die kognitive Kapazität der einzelnen Berater:in.
    \item \textbf{EBF} (Evidence-Based Framework for Economic and Social Behavior): Die vollständige wissenschaftliche Formalisierung und Systematisierung des Wissens, das hinter dem BCM steht -- und weit darüber hinaus. Über 2'000 Studien, 153 Theorien, 852 dokumentierte Fälle, formale Axiome. Das EBF \emph{enthält} das BCM, erweitert es aber massiv um theoretische Grundlagen, empirische Evidenz und methodische Werkzeuge.
    \item \textbf{BEATRIX} (Behavioral Economics AI Technology for Research and Implementation eXcellence): Verwendet digitale Technologien und maschinelles Lernen, um das EBF (und damit auch das BCM) für Berater:innen und Entscheidungsträger:innen operativ einsetzbar zu machen. BEATRIX ersetzt nicht die Beratungskompetenz, sondern hebt deren zentrale Begrenzung auf: die menschliche kognitive Kapazität.
\end{itemize}

\noindent\textit{Zusammenhang: BCM $\subset$ EBF $\leftarrow$ BEATRIX. Das BCM ist, was Berater:innen bisher im Kopf hatten. Das EBF ist alles, was dahintersteht -- formalisiert und erweitert. BEATRIX macht das EBF für Berater:innen nutzbar. Der Innovation Cheque validiert BEATRIX.}

\vspace{0.5cm}

BEATRIX ist nicht thematisch begrenzt, sondern kann grundsätzlich auf jedes Entscheidungs-, Prognose- oder Beratungsproblem angewendet werden, bei dem menschliches Verhalten eine zentrale Rolle spielt. Dies umfasst unter anderem strategische, organisatorische, marketingbezogene, anreizbezogene oder transformationsbezogene Entscheidungen. Der breite Anwendungsbereich ergibt sich daraus, dass all diese Entscheidungsprobleme auf denselben empirisch gut belegten Verhaltensmechanismen beruhen, die das EBF systematisiert und BEATRIX operativ nutzbar macht.

Ein zentraler Mehrwert von BEATRIX liegt in der Fähigkeit, Wechselwirkungen zu erfassen, die in der Realität zwischen psychologischen, finanziellen, kulturellen und institutionellen Faktoren stattfinden. In realen Entscheidungssituationen wirken diese Faktoren nahezu immer gleichzeitig und beeinflussen sich gegenseitig. Diese Interdependenzen konnten bisher aufgrund ihrer Komplexität nur begrenzt berücksichtigt werden. Das EBF modelliert diese Wechselwirkungen formal; BEATRIX macht sie erstmals konsistent analysierbar und schlägt auf dieser Basis geeignete Massnahmen vor.

Auf dieser Grundlage erlaubt BEATRIX einen systematischen Vergleich alternativer Entscheidungsverläufe. Es kann transparent analysiert werden, welche Effekte ohne Massnahmen, mit einzelnen Massnahmen oder mit unterschiedlichen Kombinationen von Massnahmen unter realistischen Kontextbedingungen zu erwarten sind. Dadurch wird insbesondere der zusätzliche Nutzen weiterer Interventionen deutlich präziser abschätzbar als mit herkömmlichen Ansätzen. Fehlpriorisierungen, ineffizienter Ressourceneinsatz und unbeabsichtigte Nebenwirkungen können so reduziert werden.

Die Innovation von BEATRIX liegt nicht in der isolierten Anwendung auf spezielle Entscheidungsprobleme, sondern in der systematischen Nutzbarmachung des im EBF formalisierten, empirisch fundierten Entscheidungswissens unter realen organisationalen Bedingungen. Der Innosuisse Innovation Cheque dient der gezielten Validierung zweier verbleibender Unsicherheiten: der technischen Skalierbarkeit dieser Entscheidungsarchitektur sowie der grundsätzlichen Markttauglichkeit und Glaubwürdigkeit einer solchen Entscheidungslogik im Beratungskontext. Die unabhängige Validierung und Weiterentwicklung von BEATRIX durch die Universität Zürich reduziert diese Risiken gezielt und schafft eine belastbare Entscheidungsgrundlage für ein mögliches nachgelagertes Innosuisse-Hauptprojekt.

\newpage

%% =============================================================================
%% 2. THEORETISCHE UND EMPIRISCHE GRUNDLAGE
%% =============================================================================

\section{Theoretische und empirische Grundlage: Vom BCM zum EBF}

Beinahe jedes unternehmerische, organisationale oder private Entscheidungsproblem zielt letztlich darauf ab, zukünftiges menschliches Verhalten zu beeinflussen oder vorherzusagen. Strategische Entscheide sollen das Verhalten von Mitarbeitenden verändern, Marketingmassnahmen das Verhalten von Konsumentinnen und Konsumenten, politische Massnahmen das Verhalten von Staatsbürgern, organisatorische Interventionen das Verhalten von Führungskräften, Teams oder Anspruchsgruppen. Unabhängig vom Anwendungskontext steht somit stets dieselbe Grundfrage im Zentrum: Wie reagieren Menschen auf veränderte Anreize, Regeln, Informationen und soziale Kontexte?

\subsection{Das BCM: Drei Jahrzehnte Beratungserfahrung}

Aus diesem Grund beruht die Beratungspraxis von FehrAdvice auf dem \textbf{Behavioral Change Model (BCM)}. Das BCM wurde über mehr als drei Jahrzehnte in enger Verbindung mit der Forschung von Prof. Ernst Fehr an der Universität Zürich und der praktischen Beratungstätigkeit von FehrAdvice entwickelt und verdichtet. Es handelt sich nicht um ein abstraktes Theoriegebäude, sondern um ein in der kontinuierlichen Auseinandersetzung mit realen Entscheidungsproblemen in Unternehmen, öffentlichen Organisationen und institutionellen Kontexten gewachsenes Beratungsmodell.

Zentral für das BCM ist die systematische Berücksichtigung von Kontextabhängigkeit menschlichen Verhaltens. Menschen legen ihrem Verhalten und ihren Urteilen nicht in allen Situationen dieselben Massstäbe zugrunde. Dieselbe Person bewertet Anreize, Fairness, Risiken oder Verpflichtungen anders, je nachdem ob sie als Arbeitnehmerin, als Konsument, als Führungskraft, als Vereinsmitglied oder als Elternteil handelt. Diese kontextabhängigen Bewertungs- und Entscheidungsregeln sind empirisch gut belegt, werden in der Praxis jedoch häufig implizit oder gar nicht berücksichtigt.

Darüber hinaus erfasst das BCM eine Vielzahl empirisch belegter Verhaltensneigungen sowie deren Heterogenität in relevanten Populationen. Menschen unterscheiden sich systematisch in Präferenzen, sozialen Motiven, Risikoaversion, Reziprozität oder Zeitpräferenzen. Diese Unterschiede sind für die Wirkung von Massnahmen zentral, können jedoch aufgrund ihrer Komplexität von einzelnen Berater:innen nicht vollständig und konsistent berücksichtigt werden.

Das BCM war mächtig -- aber begrenzt durch die kognitive Kapazität der einzelnen Berater:in. Jede Beraterin, jeder Berater konnte nur einen Teil des verfügbaren Wissens gleichzeitig abrufen und in einer konkreten Entscheidungssituation anwenden.

\subsection{Das EBF: Wissenschaftliche Formalisierung und Erweiterung}

Das \textbf{Evidence-Based Framework (EBF)} ist die vollständige wissenschaftliche Formalisierung und Systematisierung des Wissens, das hinter dem BCM steht -- und geht weit darüber hinaus. Während das BCM als mentales Rahmenwerk in den Köpfen der Berater:innen lebte, macht das EBF dieses Wissen explizit, formal und maschinenlesbar.

Das EBF umfasst über 2'000 wissenschaftliche Studien, 153 formalisierte Theorien, 852 dokumentierte Anwendungsfälle und 414 formale Axiome, systematisiert in 21 Kapiteln und 56 Appendices. Es formalisiert insbesondere drei Kernelemente, die über das BCM hinausgehen:

\begin{enumerate}
    \item \textbf{Systematische Kontextmodellierung ($\Psi$):} Acht Kontext-Dimensionen -- zeitlich, sozial, institutionell, materiell, kognitiv, ökonomisch, kulturell und physisch -- werden formal als $\Psi$-Vektor erfasst und in Entscheidungsanalysen integriert.
    \item \textbf{Formalisierte Wechselwirkungen ($\gamma$):} Die empirisch belegten Wechselwirkungen, die in der Realität zwischen psychologischen, finanziellen und kulturellen Faktoren stattfinden, werden als Komplementaritätsmatrix quantifiziert.
    \item \textbf{Kontextabhängige Parametertransformation:} Verhaltensparameter sind keine Konstanten, sondern Funktionen des Kontexts: $\theta = f(\Psi)$. Ein Anreizsystem, das in Skandinavien wirkt, kann in der Schweiz anders wirken -- das EBF modelliert systematisch, warum und wie stark.
\end{enumerate}

Das EBF enthält das BCM (BCM $\subset$ EBF), macht aber dessen implizites Wissen explizit und erweitert es um formale Grundlagen, die eine konsistente, skalierbare Anwendung ermöglichen.

\subsection{BEATRIX: Operative Nutzbarmachung des EBF}

BEATRIX nutzt das EBF nicht als abstrakten Referenzrahmen, sondern als inhaltliche Entscheidungsgrundlage, die durch KI erstmals in ihrer gesamten Breite nutzbar gemacht wird. Die KI-gestützte Architektur von BEATRIX erlaubt es, Kontexteinflüsse, Wechselwirkungen zwischen psychologischen, finanziellen und kulturellen Faktoren sowie deren Heterogenität in der relevanten Population \emph{gleichzeitig} zu erfassen und konsistent zu verarbeiten. Diese gleichzeitige Erfassung übersteigt die kognitiven Möglichkeiten einzelner Menschen oder Teams, ist jedoch durch die im EBF formalisierte und durch BEATRIX operationalisierte Entscheidungslogik möglich.

BEATRIX stellt auch die theoretische Rückführbarkeit von Analysen und Handlungsempfehlungen sicher. Dies bedeutet, dass jede Handlungsempfehlung und jede Prognose auf expliziten, empirisch gestützten theoretischen Annahmen und stabilen empirischen Regelmässigkeiten der verhaltensökonomischen Forschung beruht. Diese Annahmen sind nicht implizit oder verborgen, sondern jederzeit nachvollziehbar abrufbar. Damit ist transparent, auf welchen verhaltensökonomischen Mechanismen eine Empfehlung beruht und warum unter gegebenen Kontextbedingungen bestimmte Wirkungen erwartet werden.

Weder das BCM noch das EBF sind Gegenstand der im Innovation Cheque vorgesehenen Prüfung. Sie bilden die stabile inhaltliche Grundlage, auf der BEATRIX aufsetzt. Die beabsichtigte wissenschaftliche Validierung fokussiert sich auf die technische Skalierbarkeit von BEATRIX sowie auf die grundsätzliche Markttauglichkeit der transparent hergeleiteten, theoriegeleiteten Empfehlungen im Beratungskontext.

\newpage

%% =============================================================================
%% 3. DIE VORTEILE VON BEATRIX GEGENÜBER DER BISHERIGEN PRAXIS
%% =============================================================================

\section{Die Vorteile von BEATRIX gegenüber der bisherigen Praxis}

In komplexen organisationalen Entscheidungssituationen werden Entscheidungen heute überwiegend auf Basis individueller Erfahrung, Intuition, bekannter Heuristiken und punktueller Analysen getroffen. Auch wenn verhaltensökonomische Einsichten zunehmend Eingang in Entscheidungsprozesse finden, geschieht dies in der Regel selektiv, implizit und nicht systematisch. Empirisch belegte Verhaltensneigungen werden häufig einzeln berücksichtigt, während die Wechselwirkungen, die in der Realität zwischen psychologischen, finanziellen und kulturellen Faktoren stattfinden, die Kontextabhängigkeit dieser Faktoren sowie ihre Heterogenität in der relevanten Population nur unvollständig erfasst werden.

In realen Entscheidungssituationen wirken diese Faktoren jedoch nahezu immer gleichzeitig, unter wechselnden Kontextbedingungen und auf Populationen mit unterschiedlichen Verhaltensmustern. Ohne ein geeignetes Analysesystem, das diese Faktoren gleichzeitig und konsistent erfasst, bleibt unklar, wie sie sich gegenseitig verstärken, neutralisieren oder unbeabsichtigte Effekte erzeugen. Insbesondere können alternative Entscheidungsverläufe -- etwa das bewusste Weglassen einzelner Massnahmen oder unterschiedliche Kombinationen von Interventionen -- nicht belastbar und vergleichbar analysiert werden.

\subsection{Strukturelle Grenzen menschlicher Entscheidungs- und Beratungspraxis}

Diese Einschränkungen sind keine Frage mangelnder Kompetenz. Beratende Personen und Entscheidungsträger:innen verfügen über hohe fachliche Expertise und umfangreiche Erfahrung. Die Grenzen entstehen vielmehr aus einer strukturellen Kapazitätsbegrenzung menschlicher Entscheidungsarbeit. Die gleichzeitige Berücksichtigung mehrerer Kontexte, zahlreicher empirisch belegter Verhaltensneigungen, ihrer Heterogenität in der Population sowie der realen Wechselwirkungen zwischen psychologischen, finanziellen und kulturellen Faktoren übersteigt selbst die Möglichkeiten hochqualifizierter Individuen und Teams.

In der Praxis führt dies zwangsläufig zu Vereinfachungen. Bestimmte Verhaltensmechanismen werden explizit berücksichtigt, andere implizit ausgeblendet; Kontextabhängigkeiten werden angenähert; Wechselwirkungen nur grob eingeschätzt. Diese Vereinfachungen sind unvermeidbar, bleiben jedoch meist nicht transparent und nicht systematisch überprüfbar.

Dies war auch die zentrale Begrenzung des BCM: Solange das Wissen nur in den Köpfen der Berater:innen existierte, konnte es nie vollständig, konsistent und gleichzeitig abgerufen werden. Die Formalisierung im EBF machte das Wissen explizit; BEATRIX macht es operativ nutzbar.

\subsection{Die verbleibende strukturelle Lücke}

Genau an dieser Stelle zeigt sich die strukturelle Grenze bestehender Ansätze: Es fehlt ein Entscheidungssystem, das verhaltensökonomisches Wissen nicht nur punktuell, sondern in seiner Gesamtheit -- inklusive Kontextabhängigkeit, Heterogenität und den in der Realität stattfindenden Wechselwirkungen -- konsistent verarbeiten kann. Insbesondere fehlt die Möglichkeit, auf dieser Grundlage systematische Counterfactuals zu erzeugen und alternative Entscheidungsverläufe unter realistischen Kontextbedingungen vergleichbar zu machen.

Ein weiteres, bislang ungelöstes strukturelles Problem besteht in der fehlenden Konsistenz von Entscheidungssystemen unter Variation einzelner Annahmen oder Massnahmen. In bestehenden Ansätzen führen Änderungen an einer Kontextannahme, einer Verhaltensannahme oder einer Massnahme häufig zu lokalen Anpassungen, deren systemweite Konsequenzen nicht vollständig nachvollzogen werden. Entscheidungsmodelle verlieren dadurch bei wachsender Komplexität an Kohärenz und Stabilität. Es fehlt ein Entscheidungssystem, in dem Änderungen systemweit, konsistent und ohne ad-hoc-Korrekturen verarbeitet werden. Genau diese Fähigkeit -- die Aufrechterhaltung eines stabilen, in sich konsistenten Systems auch bei hoher Komplexität und zahlreichen Alternativszenarien -- ist mit herkömmlichen, sequenziellen oder modularen Ansätzen nicht erreichbar und bildet eine zentrale Voraussetzung für belastbare Counterfactuals.

Das EBF stellt die formale Grundlage bereit, um diese strukturelle Lücke zu schliessen. BEATRIX operationalisiert das EBF und adressiert damit genau diese Lücke: Es ermöglicht erstmals, das gesamte im EBF formalisierte verhaltensökonomische Wissen gleichzeitig abzurufen, die vielfältigen Wechselwirkungen zwischen Kontext, heterogenen Verhaltensneigungen und institutionellen Rahmenbedingungen zu erfassen -- und auf dieser Basis geeignete Massnahmen vorzuschlagen.

\newpage

%% =============================================================================
%% 4. MARKTRELEVANZ UND ORGANISATIONALER NUTZEN
%% =============================================================================

\section{Marktrelevanz und organisationaler Nutzen}

\subsection{Wo entstehen konkrete, vermeidbare Kosten in Organisationen?}

In komplexen Organisationen entstehen erhebliche Kosten nicht nur durch einzelne Fehlentscheide, sondern durch unzureichend fundierte Priorisierungen über mehrere gleichzeitig wirksame Massnahmen hinweg. Entscheidungen müssen häufig unter Zeitdruck und mit unvollständiger Transparenz über Wechselwirkungen getroffen werden. Alternative Handlungsoptionen können zwar diskutiert, aber nicht systematisch und vergleichbar bewertet werden. Die Folge sind verzögerte Anpassungen, inkonsistente Massnahmenbündel und ein Ressourceneinsatz, der sich erst im Nachhinein als suboptimal erweist. Solche Effekte betreffen nicht einzelne Sonderfälle, sondern treten regelmässig dort auf, wo mehrere Interventionen parallel umgesetzt werden und ihre Wirkungen schwer trennbar sind.

\subsection{Für welche Organisationen ist dieser Nutzen relevant?}

BEATRIX richtet sich an Organisationen, die wiederholt Entscheidungen unter hoher Komplexität treffen müssen. Dies betrifft Unternehmen, Beratungen und öffentliche Organisationen gleichermassen, insbesondere in Kontexten wie Strategieentwicklung, Organisationsgestaltung, Anreizsystemen oder Transformationsvorhaben.

In diesen Situationen steigen sowohl die Anzahl relevanter Entscheidungsdimensionen als auch die Abhängigkeiten zwischen Einflussfaktoren. Gleichzeitig sinkt die Fähigkeit, unterschiedliche Handlungsoptionen konsistent zu vergleichen. Genau hier entsteht der Bedarf nach einer Entscheidungsgrundlage, die über intuitive Abwägungen und isolierte Analysen hinausgeht.

\subsection{Welcher konkrete Nutzen entsteht durch BEATRIX?}

BEATRIX verändert nicht nur einzelne Entscheidungsinhalte, sondern auch die Qualität der Entscheidungsgrundlage. Organisationen erhalten die Möglichkeit, Handlungsoptionen systematisch zu vergleichen, bevor sie umgesetzt werden. Dies umfasst nicht nur die Bewertung einzelner Massnahmen, sondern insbesondere die Erfassung der in der Realität stattfindenden Wechselwirkungen zwischen psychologischen, finanziellen und kulturellen Faktoren -- sowie die explizite Berücksichtigung von Alternativen, einschliesslich der Option, auf bestimmte Massnahmen bewusst zu verzichten.

Dadurch werden Entscheidungen früher überprüfbar und Prioritäten klarer begründbar. Ressourcen können gezielter eingesetzt, Massnahmen konsistenter aufeinander abgestimmt und unerwünschte Nebenwirkungen frühzeitig erkannt werden. Der Nutzen entsteht somit sowohl durch die einzelnen, weit besser fundierten Empfehlungen als auch durch eine höhere Robustheit von Entscheidungsprozessen unter komplexen Bedingungen.

\subsection{Anwendungsbeispiel 1: Regulierte energiepolitische Kontexte}

In der Schweiz stehen energiepolitische Entscheide regelmässig vor der Herausforderung, mehrere parallel wirksame Massnahmen unter komplexen regulatorischen, föderalen und gesellschaftlichen Rahmenbedingungen abzustimmen. Dazu gehören beispielsweise Förderinstrumente, Informationskampagnen, regulatorische Vorgaben oder finanzielle Anreizsysteme, deren Wirkungen sich gegenseitig beeinflussen.

BEATRIX kann in solchen Kontexten eingesetzt werden, um alternative Massnahmenbündel systematisch zu vergleichen. Unter Berücksichtigung unterschiedlicher Akteursgruppen, institutioneller Zuständigkeiten und relevanter Kontextbedingungen erlaubt BEATRIX eine strukturierte Prognose möglicher Wirkungen, Zielkonflikte und unbeabsichtigter Nebenwirkungen, bevor Entscheidungen umgesetzt werden. Damit unterstützt BEATRIX Entscheidungsträger:innen dabei, Prioritäten nachvollziehbar zu setzen und Massnahmen konsistent aufeinander abzustimmen.

\subsection{Anwendungsbeispiel 2: Anreiz- und Organisationssysteme im Gesundheitswesen}

Im Schweizer Gesundheitswesen müssen Organisationen regelmässig Entscheidungen über Anreizsysteme, Organisationsstrukturen oder Veränderungsmassnahmen treffen, deren Wirkungen sich über mehrere Ebenen hinweg entfalten. Finanzielle Anreize, organisatorische Vorgaben, Informationsinstrumente und professionelle Normen wirken dabei gleichzeitig und können sich gegenseitig verstärken oder neutralisieren.

BEATRIX kann in solchen Entscheidungssituationen eingesetzt werden, um unterschiedliche Gestaltungsoptionen vergleichbar zu machen. Durch die gleichzeitige Berücksichtigung relevanter Kontextbedingungen und heterogener Verhaltensreaktionen ermöglicht BEATRIX eine strukturierte Einschätzung, welche Kombinationen von Massnahmen unter gegebenen Bedingungen voraussichtlich wirksam sind und wo Risiken unbeabsichtigter Effekte bestehen.

\subsection{Anwendungsbeispiel 3: Akquisitionsstrategien in kulturell heterogenen Wachstumsmärkten}

Unternehmen, die durch Akquisitionen in Wachstumsmärkten expandieren, stehen vor einer besonderen Herausforderung: Entlang der gesamten Transaktion -- von der Investitionslogik über die Prüfung potenzieller Zielunternehmen und die Verhandlungsführung bis zur Integration nach dem Closing -- wirken unterschiedliche Verhaltensfaktoren gleichzeitig und beeinflussen sich gegenseitig. Überschätzung des eigenen Integrationsvermögens auf Führungsebene, kulturelle Inkompatibilität zwischen den Organisationen, kontextabhängige Verlustaversion in der Verhandlung sowie konkurrierende Anreizsysteme in der Integrationsphase werden in bestehenden Ansätzen typischerweise getrennt nach Transaktionsphase oder Fachdisziplin analysiert. Die Wechselwirkungen, die in der Realität zwischen diesen psychologischen, finanziellen und kulturellen Faktoren über Transaktionsphasen hinweg stattfinden, bleiben damit unberücksichtigt -- insbesondere in Märkten mit grosser kultureller Distanz, wo erfahrungsbasierte Einschätzungen am wenigsten verlässlich sind.

BEATRIX ermöglicht es, eine Akquisition als konsistentes Gesamtsystem über alle Transaktionsphasen hinweg zu modellieren. Die kulturelle Kompatibilität zwischen Käufer und Zielunternehmen kann auf denselben formalen Kontextdimensionen gemessen werden wie organisatorische Faktoren oder Anreizstrukturen. Wechselwirkungen -- beispielsweise zwischen finanziellen Integrationsanreizen und der Bereitschaft lokaler Führungskräfte zur Zusammenarbeit, oder zwischen der gewählten Integrationsstrategie und der Innovationsfähigkeit des Zielunternehmens -- werden als explizite Komplementaritäten erfasst. Alternative Integrationspfade lassen sich als Counterfactuals unter variierten Kontextannahmen strukturiert vergleichen. Damit unterstützt BEATRIX Entscheidungsträger:innen dabei, Akquisitionsrisiken frühzeitig zu identifizieren, Integrationsstrategien konsistent auf die jeweiligen kulturellen und organisatorischen Bedingungen abzustimmen und den Ressourceneinsatz über den gesamten Transaktionsprozess hinweg gezielter zu steuern.

\subsection{Organisationaler Mehrwert über einzelne Entscheidungen hinaus}

Über einzelne Entscheidungsprozesse hinaus trägt BEATRIX zu einer höheren Transparenz und Nachvollziehbarkeit von Entscheidungen bei. Entscheidungsannahmen, Abwägungen und Alternativen werden explizit gemacht und sind auch ex post überprüfbar. Dies erleichtert interne Abstimmungen, erhöht die Anschlussfähigkeit von Entscheidungen über Hierarchie- und Funktionsgrenzen hinweg und unterstützt eine konsistente Entscheidungsgovernance.

Für Beratungen und interne Strategiefunktionen eröffnet BEATRIX zudem die Möglichkeit, Entscheidungslogiken systematisch zu dokumentieren und über Projekte hinweg weiterzuentwickeln. Damit entsteht ein nachhaltiger organisationaler Nutzen, der über den Einzelfall hinaus wirkt.

\subsection{Warum dieser Nutzen marktrelevant ist}

Die beschriebenen Vorteile adressieren ein strukturelles Problem moderner Organisationen: die steigende Komplexität von Entscheidungsumfeldern bei gleichzeitig begrenzter Entscheidungsfähigkeit. Dieses Problem zeigt sich in der Schweiz insbesondere in Organisationen und Märkten, die durch hohe regulatorische Anforderungen, ausgeprägte Koordinationsbedarfe zwischen Akteuren sowie politisch und gesellschaftlich sensible Rahmenbedingungen geprägt sind.

BEATRIX bietet vor diesem Hintergrund einen Effizienzgewinn durch fundierte Entscheidungen und eine Entscheidungsunterstützung, die dort ansetzt, wo bestehende Instrumente und Methoden in solchen komplexen, stark kontextabhängigen Umfeldern an ihre Grenzen stossen.

Damit adressiert BEATRIX einen klar identifizierbaren Bedarf bei Organisationen, die regelmässig mit komplexen, mehrdimensionalen Entscheidungsproblemen konfrontiert sind -- und schafft einen nachhaltigen, differenzierenden Nutzen.

\newpage

%% =============================================================================
%% 5. FEASIBILITY & RISK REDUCTION
%% =============================================================================

\section{Feasibility \& Risk Reduction}

\subsection{Was existiert bereits}

BEATRIX verfügt über einen funktionsfähigen Kern, der bereits heute in Beratungs- und Entscheidungskontexten bei FehrAdvice eingesetzt wird. Die zugrunde liegende Entscheidungslogik -- aufbauend auf dem im EBF formalisierten verhaltensökonomischen Wissen -- erlaubt es, relevante Kontextbedingungen, unterschiedliche Einflussfaktoren sowie deren Wechselwirkungen gleichzeitig und konsistent zu berücksichtigen. Diese gleichzeitige Verarbeitung ersetzt keine menschliche Entscheidungsarbeit, sondern stellt sicher, dass bei Analysen und Prognosen alle relevanten Einflussfaktoren parallel berücksichtigt werden und nicht -- wie in der bisherigen Praxis üblich -- schrittweise oder isoliert.

Der aktuelle Entwicklungsstand erlaubt damit fundierte Analysen unter realistischen Anwendungsbedingungen und zeigt, dass BEATRIX im Beratungsalltag funktioniert. Bislang liegt jedoch keine unabhängige externe Validierung vor, weder im Hinblick auf die technische Tragfähigkeit und Skalierbarkeit dieser Entscheidungsarchitektur noch im Hinblick auf das Marktpotenzial eines theoriegeleiteten Prognose- und Beratungstools dieser Art. Genau diese beiden offenen Punkte begrenzen aktuell eine weitergehende Skalierung und Markterschliessung.

Ein wesentlicher Teil des Entwicklungsrisikos ist durch das zugrunde liegende BCM und dessen Formalisierung im EBF reduziert. Das BCM stellt ein über mehrere Jahrzehnte entwickeltes, theoriegeleitetes und empirisch fundiertes Beratungsmodell dar, das sowohl wissenschaftlich breit abgestützt als auch in der praktischen Beratung erprobt ist. Das EBF erweitert dieses Beratungsmodell um die vollständige wissenschaftliche Formalisierung, systematische Evidenzbasis und formale Axiomatik. Weder das BCM noch das EBF sind Gegenstand der im Innovation Cheque vorgesehenen Prüfung. Sie bilden die stabile inhaltliche Grundlage, auf der BEATRIX aufsetzt.

Gegenstand der im Innovation Cheque vorgesehenen Validierung ist vielmehr die Frage, ob die auf dem EBF aufbauende Entscheidungslogik von BEATRIX -- erstens technisch tragfähig skaliert werden kann, wenn sie alle relevanten Einflussfaktoren gleichzeitig verarbeitet, und -- zweitens als Grundlage für Beratung und Prognose im Markt grundsätzlich anschlussfähig ist.

Dabei geht es ausdrücklich nicht um umfassende Akzeptanztests im Sinne von Feldexperimenten, breiten Nutzerstudien oder quantitativen Surveys. Solche Untersuchungen setzen eine weiterentwickelte Anwendung und reale Markteinführung voraus und sind Bestandteil eines möglichen Innosuisse-Hauptprojekts. Im Rahmen des Innovation Cheque wird hingegen geklärt, ob die Entscheidungsarchitektur und ihre transparente, theoriegeleitete Herleitung prinzipiell geeignet sind, in Beratungs- und Prognosekontexten eingesetzt und weiter skaliert zu werden.

Die Einbindung der Universität Zürich dient gezielt der unabhängigen Reduktion dieser verbleibenden Risiken. Die im Innovation Cheque adressierten Unsicherheiten betreffen ausschliesslich (i) die technische Skalierbarkeit einer Entscheidungsarchitektur, die komplexe Entscheidungssituationen ganzheitlich und gleichzeitig verarbeitet, sowie (ii) die grundsätzliche Markttauglichkeit und Glaubwürdigkeit einer solchen Entscheidungslogik im Beratungskontext.

Diese beiden Aspekte werden unabhängig untersucht: Professor Harald Gall prüft die technische Tragfähigkeit und Skalierbarkeit der Architektur, während Professor Johannes Luger empirisch analysiert, unter welchen Bedingungen transparent hergeleitete, theoriegeleitete Entscheidungslogiken für Beratung und Prognose als nutzbar und plausibel wahrgenommen werden.

Die Research Collaboration bildet damit eine zielgerichtete, anwendungsnahe Risikoabsicherung für eine mögliche Weiterentwicklung von BEATRIX in einem Innosuisse-Hauptprojekt. Ein solches Hauptprojekt würde auf realen Anwendungsszenarien, marktorientierten Fragestellungen sowie dem Einsatz moderner Plattform- und KI-Technologien aufbauen und die im Innovation Cheque identifizierten offenen Punkte systematisch weiterführen.

\subsection{Welche Risiken sind offen}

Zwei zentrale Risiken sind für den nächsten Entwicklungsschritt entscheidend: Erstens die technische Skalierbarkeit der Entscheidungsarchitektur bei wachsender Komplexität und Nutzung. Zweitens die grundsätzliche Markttauglichkeit theoriegeleiteter, transparent hergeleiteter Empfehlungen im Beratungskontext, insbesondere im Vergleich zu erfahrungsbasierten Entscheidungsansätzen. Diese Risiken lassen sich nicht intern abschliessend beurteilen und erfordern eine unabhängige wissenschaftliche Prüfung.

\subsection{Was CHF 15'000 konkret klärt}

Der Innovation Cheque ermöglicht eine fokussierte Machbarkeitsstudie, die genau diese Risiken adressiert. Mit den bereitgestellten Mitteln werden (i) eine strukturierte Analyse der technischen Architektur im Hinblick auf Skalierbarkeit sowie (ii) eine empirische Untersuchung der Markttauglichkeit mit Führungskräften durchgeführt. Die Ergebnisse liefern eine belastbare Entscheidungsgrundlage für den nächsten Entwicklungsschritt.

Ohne den Innovation Cheque könnten diese Risiken nur ex post beurteilt werden. Der Cheque reduziert damit gezielt das Risiko von Fehlentwicklungen, Fehlinvestitionen und nicht skalierbaren Ansätzen vor einem möglichen Hauptprojekt.

\subsection{Klare Go/No-Go-Kriterien}

Am Ende der Machbarkeitsstudie steht eine klare Entscheidung: Ein Go erfolgt bei nachgewiesener Skalierbarkeit und ausreichender Markttauglichkeit der Entscheidungslogik. Ein No-Go erfolgt, wenn fundamentale architektonische Grenzen oder mangelnde Anschlussfähigkeit identifiziert werden. Damit reduziert der Innovation Cheque gezielt Investitions- und Entwicklungsrisiken.

\subsection{Perspektive eines nachgelagerten Innosuisse-Hauptprojekts}

Ein erfolgreich abgeschlossener Innovation Cheque schafft die Voraussetzung, BEATRIX in einem nachgelagerten Innosuisse-Hauptprojekt gezielt weiterzuentwickeln. Ziel eines solchen Hauptprojekts wäre nicht die Weiterentwicklung des BCM oder des EBF, sondern die anwendungsnahe Ausgestaltung einer skalierbaren Entscheidungs- und Beratungsarchitektur auf Basis der bereits validierten Entscheidungslogik.

Konkret würde ein Hauptprojekt erlauben, BEATRIX über ein reines textbasiertes Interface hinaus weiterzuentwickeln und eine strukturierte, nutzerfreundliche Interaktion zu gestalten. Dazu gehört insbesondere die Frage, wie BEATRIX Nutzer:innen schrittweise durch komplexe Entscheidungs- und Prognoseprozesse führen kann, etwa indem relevante Rückfragen selbständig gestellt und erforderliche Eingaben systematisch eingefordert werden. Ziel wäre eine Anwendung, die Entscheidungsprobleme nicht nur beantwortet, sondern deren Strukturierung aktiv unterstützt.

Ein weiterer Fokus eines Hauptprojekts läge auf der Klärung unterschiedlicher Nutzungsszenarien. Dies betrifft insbesondere das Zusammenspiel zwischen BEATRIX als Entscheidungsarchitektur und der Beratungspraxis von FehrAdvice: von Anwendungen, bei denen BEATRIX primär als internes Beratungsinstrument eingesetzt wird, bis hin zu Szenarien, in denen Nutzer:innen bestimmte Fragestellungen selbständig bearbeiten können, während FehrAdvice eine begleitende oder vertiefende Rolle übernimmt.

Darüber hinaus würde ein Hauptprojekt die Möglichkeit eröffnen, BEATRIX als softwarebasierte Plattform weiterzuentwickeln, die in unterschiedlichen Beratungs- und Prognosedomänen einsetzbar ist. Dabei stünde nicht die Entwicklung einzelner Speziallösungen im Vordergrund, sondern die Frage, wie eine einheitliche Entscheidungslogik -- aufbauend auf dem EBF -- domänenspezifisch ausgestaltet werden kann, etwa über angepasste Benutzeroberflächen und Anwendungskontexte.

Der Innovation Cheque ist Voraussetzung für diese Weiterentwicklung, da zentrale Entscheidungen über Nutzerführung, Plattformarchitektur und Anwendungslogik erst dann seriös getroffen werden können, wenn die technische Skalierbarkeit der Entscheidungsarchitektur sowie ihre grundsätzliche Markttauglichkeit unabhängig validiert sind.

\newpage

%% =============================================================================
%% 6. RESEARCH COLLABORATION
%% =============================================================================

\section{Research Collaboration}

\subsection{Warum diese Kompetenzen fehlen}

Das EBF -- aufbauend auf dem über Jahrzehnte entwickelten BCM -- bildet die theoretische und empirische Grundlage von BEATRIX und ist nicht Gegenstand der wissenschaftlichen Prüfung im Rahmen des Innovation Cheque. Das Modell ist wissenschaftlich breit abgestützt und in der praktischen Beratung erprobt. Die Research Collaboration fokussiert sich daher gezielt auf zwei verbleibende, für eine Weiterentwicklung entscheidende Unsicherheiten.

Erstens betrifft dies die technische Skalierbarkeit der Entscheidungsarchitektur von BEATRIX. Konkret geht es um die Frage, ob eine Entscheidungslogik, die relevante Kontextbedingungen, Einflussfaktoren und deren Wechselwirkungen gleichzeitig und konsistent verarbeitet, auch bei wachsender Komplexität technisch tragfähig bleibt.

Zweitens adressiert die Research Collaboration die grundsätzliche Anschlussfähigkeit transparenter, theoriegeleiteter Entscheidungslogiken im Beratungskontext. Dabei geht es nicht um umfassende Akzeptanztests im Sinne von Feldexperimenten oder breiten Nutzerstudien, sondern um die empirische Klärung, unter welchen Bedingungen solche Entscheidungsgrundlagen für Beratung und Prognose als plausibel und nutzbar wahrgenommen werden.

Damit adressiert die Zusammenarbeit exakt jene Punkte, die für eine mögliche Weiterentwicklung von BEATRIX in ein Innosuisse-Hauptprojekt entscheidend sind. Beide Fragestellungen erfordern unabhängige wissenschaftliche Expertise und können weder intern noch rein technisch abschliessend beurteilt werden.

FehrAdvice verfügt über umfassende verhaltensökonomische Expertise und praktische Erfahrung in der Anwendung des EBF durch BEATRIX. Was intern nicht abgedeckt werden kann, ist die unabhängige wissenschaftliche Prüfung der technischen Skalierbarkeit einer Entscheidungsarchitektur, die das gesamte im EBF formalisierte Wissen gleichzeitig verarbeitet, sowie der empirischen Markttauglichkeit solcher theoriegeleiteter Empfehlungen bei Führungskräften. Diese Kompetenzen erfordern methodische Tiefe und wissenschaftliche Neutralität, die ausserhalb der Organisation liegen müssen.

\subsection{Warum die Universität Zürich}

Die Universität Zürich vereint genau jene wissenschaftlichen Kompetenzen, die für die unabhängige Validierung von BEATRIX im Rahmen des Innovation Cheque erforderlich sind. Sie verbindet verhaltensökonomische Grundlagenforschung mit evidenzbasierter Strategieforschung sowie ausgewiesener Expertise in Software-, Plattform- und KI-Architekturen. Diese Kombination ist notwendig, um sowohl die inhaltliche Entscheidungslogik als auch deren technische Tragfähigkeit und praktische Anschlussfähigkeit fundiert beurteilen zu können.

Im Zusammenspiel mit der anwendungsorientierten Beratungspraxis von FehrAdvice entsteht dadurch ein klar abgegrenzter Mehrwert: Während FehrAdvice BEATRIX im Beratungsalltag einsetzt und weiterentwickelt, übernimmt die Universität Zürich die unabhängige wissenschaftliche Prüfung der zugrunde liegenden Entscheidungsarchitektur. Gegenstand der Zusammenarbeit ist damit nicht die gemeinsame Produktentwicklung, sondern die Validierung der technischen Skalierbarkeit und der Markttauglichkeit einer auf dem EBF aufbauenden, theoriegeleiteten Entscheidungslogik.

Die institutionelle Einbettung an der Universität Zürich stellt dabei wissenschaftliche Unabhängigkeit, methodische Qualität und Glaubwürdigkeit der Ergebnisse sicher. Die räumliche Nähe sowie bestehende Kooperationsstrukturen ermöglichen eine effiziente und zielgerichtete Durchführung der Machbarkeitsstudie im Rahmen des Innovation Cheque.

\subsection{Warum Prof. Gall \& Prof. Luger}

\textbf{Prof. Harald Gall} (Institut für Informatik, Dekan der Wirtschaftswissenschaftlichen Fakultät der UZH)

Renommierter Experte für Software-Entwicklung mit über 300 Publikationen. In Forschung und Lehre beschäftigt er sich mit Softwarearchitektur und -design, Softwarequalität, KI-gestützter Softwareentwicklung, skalierbaren Daten- und Softwarelösungen sowie dem Transfer von Innovationen in das Unternehmertum. Im BEATRIX-Projekt verantwortet er die Entwicklung skalierbarer Plattformansätze sowie das domänen- und marktspezifische Design für den Einsatz von BEATRIX. Seine Frage:

\textit{``Wie muss ein System wie BEATRIX gebaut sein, damit es zuverlässig wächst und sich weiterentwickelt?''}

\textbf{Prof. Johannes Luger} (Institut für Betriebswirtschaft)

Forscht zur Zusammenarbeit zwischen Menschen und intelligenten Systemen (Krakowski et al., 2025) mit ausgewiesener Expertise in der empirischen Untersuchung von Entscheidungsverhalten im Top-Management. Seine Forschung fokussiert sich darauf, unter welchen Bedingungen Führungskräfte transparent hergeleiteten, theoriegeleiteten Entscheidungslogiken Vertrauen schenken und diese tatsächlich nutzen. Im BEATRIX-Projekt prüft er empirisch die Markttauglichkeit der auf dem EBF aufbauenden Entscheidungsarchitektur bei Führungskräften. Seine Frage:

\textit{``Unter welchen Bedingungen vertrauen Führungskräfte den Empfehlungen eines theoriegeleiteten Analyse-Systems?''}

\subsection{Anschlussfähigkeit an ein Hauptprojekt}

Die Ergebnisse der Machbarkeitsstudie liefern eine belastbare Grundlage für ein Innosuisse-Hauptprojekt. Sie klären, ob und wie BEATRIX technisch skaliert und im Beratungsmarkt eingeführt werden kann. Damit schafft die Zusammenarbeit eine klare Entscheidungsbasis für den nächsten Entwicklungsschritt.

Die Universität Zürich entwickelt BEATRIX nicht mit, sondern validiert die bestehende Entscheidungsarchitektur unabhängig und wissenschaftlich.

\newpage

%% =============================================================================
%% 7. IMPACT FOR SWITZERLAND
%% =============================================================================

\section{Impact for Switzerland}

\subsection{IP in der Schweiz}

BEATRIX basiert auf einer wissenschaftlich fundierten Entscheidungsarchitektur -- dem EBF, aufbauend auf dem über Jahrzehnte entwickelten BCM --, deren geistiges Eigentum vollständig in der Schweiz entwickelt, weiterentwickelt und verankert ist. Die im Rahmen des Innovation Cheque vorgesehene Validierung stärkt die Position der Schweiz nicht auf Produktebene, sondern auf Ebene methodischer Kernkompetenzen: der Entwicklung, Absicherung und praktischen Nutzbarmachung evidenzbasierter Entscheidungsarchitekturen für komplexe organisationale Kontexte.

Ein zentraler Impact für die Schweiz liegt in den wirtschaftlichen Verwertungsmöglichkeiten dieser Entscheidungsarchitektur. Das in der Schweiz verankerte IP kann in unterschiedlichen Formen genutzt werden, etwa durch lizenzbasierte Nutzung von BEATRIX in Beratungs-, Organisations- oder Prognosekontexten, durch projektbasierte Anwendungen in regulierten Märkten oder durch den Aufbau skalierbarer, plattformbasierter Angebote. Damit verbleibt die Wertschöpfung aus der Weiterentwicklung und Anwendung dieser Entscheidungslogik langfristig im Schweizer Innovations- und Wirtschaftsraum.

Das Projekt adressiert damit eine strategisch relevante Fähigkeit: die Entwicklung von Entscheidungsinstrumenten, die auch unter steigender Komplexität zuverlässig einsetzbar sind und sowohl für private Unternehmen als auch für öffentliche und halböffentliche Organisationen institutionell anschlussfähig bleiben. Durch die unabhängige Validierung im Rahmen des Innovation Cheque wird eine belastbare Grundlage geschaffen, um diese Entscheidungsarchitektur gezielt weiterzuentwickeln und wirtschaftlich zu verwerten.

Der Impact für die Schweiz besteht somit nicht nur in einem einzelnen Anwendungsfall, sondern im Aufbau und der Absicherung einer hochwertigen, in der Schweiz verankerten Entscheidungs- und Beratungskompetenz, die langfristig skalierbar ist und Anschluss an weiterführende Innovations- und Marktprojekte erlaubt.

\subsection{Know-how-Aufbau}

Das Projekt fördert gezielt den Transfer verhaltensökonomischer Forschung in konkret einsetzbare Entscheidungsarchitekturen, die in Form einer softwarebasierten Plattform für Beratungs- und Prognoseanwendungen nutzbar gemacht werden. Im Zentrum steht dabei nicht die Entwicklung eines einzelnen Produkts, sondern der systematische Aufbau von Know-how darüber, wie empirisch fundierte Entscheidungslogiken -- formalisiert im EBF -- so strukturiert, implementiert und betrieben werden können, dass sie in unterschiedlichen Anwendungsdomänen replizierbar und skalierbar einsetzbar sind.

An der Schnittstelle von Sozialwissenschaften, Management und Technologie entsteht damit eine Kompetenz, die über den konkreten Anwendungsfall hinaus relevant ist: Wie lassen sich verhaltensökonomische Entscheidungsmodelle in einer Software-Architektur abbilden, die eine domänenspezifische Ausgestaltung -- etwa über angepasste Benutzeroberflächen und Anwendungskontexte in unterschiedlichen Beratungs- und Prognosebereichen -- ermöglicht, ohne die zugrunde liegende Entscheidungslogik jeweils neu zu entwickeln.

Die Zusammenarbeit zwischen FehrAdvice und der Universität Zürich trägt dazu bei, diese Kompetenz langfristig in der Schweiz zu verankern. Während FehrAdvice die anwendungsnahe Weiterentwicklung und den Einsatz solcher Entscheidungsarchitekturen im Beratungskontext vorantreibt, leistet die Universität Zürich einen Beitrag zur methodischen Fundierung und Validierung der zugrunde liegenden Konzepte. Dadurch entsteht in der Schweiz Know-how im Bereich evidenzbasiertes Management, das sowohl wissenschaftlich anschlussfähig als auch marktorientiert nutzbar ist und Perspektiven für skalierbare Anwendungen in verschiedenen Domänen eröffnet.

\subsection{Vorbildcharakter für KMU--Hochschul-Zusammenarbeit}

BEATRIX steht exemplarisch für eine Form der Zusammenarbeit zwischen KMU und Hochschule, die über klassischen Wissenstransfer hinausgeht. Im Zentrum steht nicht die Entwicklung eines einzelnen Produkts, sondern die gemeinsame Validierung einer methodischen Innovation mit klar definierten Go/No-Go-Kriterien.

Dieses Vorgehen ist übertragbar auf andere Innovationsvorhaben, bei denen wissenschaftliche Erkenntnisse in institutionell belastbare Entscheidungsinstrumente überführt werden sollen. Damit leistet das Projekt einen Beitrag zur Stärkung eines Innovationsmodells, das fokussiert, risikobewusst und wirkungsorientiert ist -- im direkten Sinne des Innosuisse-Mandats.

Die Ergebnisse der Machbarkeitsstudie erlauben eine klare, evidenzbasierte Entscheidung über die Einreichung eines Innosuisse-Hauptprojekts.

\vspace{2cm}

\begin{center}
\rule{0.4\textwidth}{0.5pt}\\[0.5em]
{\small Dokument erstellt: 27. Januar 2026 | Aktualisiert: 14. Februar 2026}\\
{\small Version: 5.2 -- Originaldokument mit BCM/EBF/BEATRIX-Abgrenzung}\\
{\small Für: Innosuisse Innovation Cheque (SSBM)}
\end{center}

\end{document}
